%! Unit 1: Functions and Change — Summative Assessment — Practice Test --- Study Copy (blank)
%! Published by `make -C unit01/tests all` to unit01/sample_test/main.pdf, which
%! shared/unit.mk merges at the tail of the unit STUDENT packet. Its key
%! (test_keys/practice_test_key) is published to sample_test_key and merged at
%! the tail of the unit KEY packet, so the two must paginate identically:
%! every worked solution lives in a \begin{work} block authored byte-identically
%! in both files. Scoring guidance goes on page 2 of unit01/unit_cover_key/,
%! never here.
\documentclass[10pt]{article}
\usepackage{precalculus-article}
\usepackage{precalculus-boxes}

% Part-divider strip
\newcommand{\parthead}[1]{%
  \vspace{6pt}\noindent
  \begin{headlinebox}{plum}{\color{white}\bfseries #1}\end{headlinebox}%
  \vspace{4pt}}

\begin{document}

\pageheader{Unit 1: Functions and Change}{Practice Test --- Study Copy}
\namedateperiod

\vspace{0.08in}
\begin{remindbox}
\small
\textbf{This is a practice test.} It mirrors the real Unit 1 test in length,
format, and the ideas it covers --- the numbers and the contexts are the only
things that differ. Work every problem with no notes, then check yourself
against the key. \textbf{80 points total.}
\end{remindbox}

% ======================================================================
\boxguard[6]
\parthead{Part A --- Vocabulary (10 pts)}
Write the letter of the best description next to each term.

\vspace{4pt}
\noindent
\begin{minipage}[t]{0.36\textwidth}
\begin{enumerate}[leftmargin=2.2em, itemsep=6pt, label=\arabic*.]
  \item Function \blank{1.1cm}
  \item Domain \blank{1.1cm}
  \item Range \blank{1.1cm}
  \item Zero of a function \blank{1.1cm}
  \item Concave up \blank{1.1cm}
  \item Average rate of change \blank{1.1cm}
  \item Parent function \blank{1.1cm}
  \item Composition \blank{1.1cm}
  \item One-to-one \blank{1.1cm}
  \item Inverse function \blank{1.1cm}
\end{enumerate}
\end{minipage}%
\hfill
\begin{minipage}[t]{0.61\textwidth}
\small
\begin{description}[leftmargin=1.4em, itemsep=2pt]
  \item[(A)] The set of every output the rule actually produces.
  \item[(B)] Feeding the output of one rule into another: $f(g(x))$.
  \item[(C)] A rule that pairs each input with exactly one output.
  \item[(D)] The rule that gives back the question: if $f(a) = b$, it sends $b$ to $a$.
  \item[(E)] An input whose output is $0$.
  \item[(F)] The simplest member of a family, before any shift or stretch.
  \item[(G)] The set of every input the rule is allowed to take.
  \item[(H)] The one steady rate that would give the same change in output over the same run.
  \item[(I)] Bending upward --- the rate of change is increasing.
  \item[(J)] No two different inputs share an output.
\end{description}
\end{minipage}

% ======================================================================
\boxguard[6]
\parthead{Part B --- Multiple Choice (12 pts)}
Circle the best answer, then write its letter in the blank at the right.

\begin{enumerate}[leftmargin=1.9em, itemsep=8pt]
  \item The statement $f(3) = 7$ tells you that
    \\[2pt]\hspace*{1em}
    (A) $f(7) = 3$ \quad (B) $7$ is an input of $f$ \quad
    (C) $(3,7)$ is on the graph of $f$ \quad (D) $f$ is increasing at $x = 3$
    \hfill \blank{1.1cm}
  \item On an interval where $f$ is decreasing and concave up, the outputs are
    \\[2pt]\hspace*{1em}
    (A) falling faster and faster \quad (B) rising more and more slowly \quad
    (C) changing at a constant rate \quad (D) falling, but more and more slowly
    \hfill \blank{1.1cm}
  \item A table of $f$ has equally spaced inputs and \textbf{constant second}
    differences. The best model is
    \\[2pt]\hspace*{1em}
    (A) linear \quad (B) quadratic \quad (C) exponential \quad (D) none of these
    \hfill \blank{1.1cm}
  \item The graph of $g(x) = f(x + 3)$ is the graph of $f$ shifted
    \\[2pt]\hspace*{1em}
    (A) $3$ left \quad (B) $3$ right \quad (C) $3$ up \quad (D) $3$ down
    \hfill \blank{1.1cm}
  \item If $f(x) = x^2$ and $g(x) = x - 1$, then $(f \circ g)(4)$ equals
    \\[2pt]\hspace*{1em}
    (A) $15$ \quad (B) $16$ \quad (C) $9$ \quad (D) $3$
    \hfill \blank{1.1cm}
  \item A function has an inverse that is also a function exactly when the
    function is
    \\[2pt]\hspace*{1em}
    (A) increasing somewhere \quad (B) concave up \quad
    (C) defined at $0$ \quad (D) one-to-one
    \hfill \blank{1.1cm}
\end{enumerate}

% ======================================================================
\boxguard[6]
\parthead{Part C --- Short Answer \& Computation (40 pts, 5 each)}
Show your work in the space provided.

\begin{enumerate}[leftmargin=1.9em, itemsep=10pt]

  \item Let $f(x) = 2x^2 - 5$.

\begin{work}
     f(3) &= 2(3)^2 - 5 \\
          &= 13 \\
    f(-2) &= 2(-2)^2 - 5 \\
          &= 3 \\
 2x^2 - 5 &= 3 \\
      x^2 &= 4 \\
        x &= \pm 2
\end{work}

    \begin{enumerate}[label=(\alph*), itemsep=3pt]
      \item $f(3) = $ \blank{1.6cm} \qquad \textbf{(b)} $f(-2) = $ \blank{1.6cm}
      \setcounter{enumii}{2}
      \item Solve $f(x) = 3$. There are two answers: \blank{2.6cm}
      \item Which of (a)--(c) is a \emph{solve} question rather than an
            \emph{evaluate} question? \blank{1.6cm}
    \end{enumerate}

  \item The graph of $y = f(x)$ is shown below.

        \vspace{2pt}
        \begin{center}
        \begin{tikzpicture}[scale=0.62]
          \draw[linegray, very thin] (-0.4,-2.4) grid (6.4,3.4);
          \draw[->] (-0.4,0) -- (6.6,0) node[right]{\small $x$};
          \draw[->] (0,-2.4) -- (0,3.6) node[above]{\small $y$};
          \foreach \x in {1,...,6} \draw (\x,0.09) -- (\x,-0.09) node[below]{\scriptsize $\x$};
          \foreach \y in {-2,-1,1,2,3} \draw (0.09,\y) -- (-0.09,\y) node[left]{\scriptsize $\y$};
          \draw[plum, very thick] plot[smooth] coordinates
            {(0,3) (1,0) (2,-1.5) (3,-2) (4,-1.5) (5,0) (6,3)};
        \end{tikzpicture}
        \end{center}

    \begin{enumerate}[label=(\alph*), itemsep=3pt]
      \item On what interval is $f$ decreasing? \blank{3.2cm}
      \item Concave up or concave down on the whole picture? \blank{3.2cm}
      \item List the zeros of $f$: \blank{3.2cm}
      \item $f(3) = $ \blank{1.6cm} \qquad \textbf{(e)} Range on the window
            shown: \blank{3cm}
    \end{enumerate}

  \item A tank is draining. Its volume $V$ (gallons) after $t$ minutes:

        \vspace{2pt}
        \begin{center}
        \renewcommand{\arraystretch}{1.3}
        \begin{tabular}{|c|c|c|c|c|}
        \hline
        $t$ (min) & $0$ & $4$ & $8$ & $12$ \\ \hline
        $V$ (gal) & $180$ & $148$ & $124$ & $108$ \\ \hline
        \end{tabular}
        \end{center}

\begin{work}
  \frac{V(8) - V(0)}{8 - 0} &= \frac{124 - 180}{8} \\
                            &= -7
\end{work}

    \begin{enumerate}[label=(\alph*), itemsep=3pt]
      \item Average rate of change of $V$ from $t = 0$ to $t = 8$:
            \blank{2cm}, with units \blank{4cm}.
      \item In one sentence, what does that number mean for the tank?

            \writeline
      \item Is the tank draining faster or slower as time goes on?
            \blank{1.8cm} \ Justify with the numbers.

            \writeline
    \end{enumerate}

  \item The table below has equally spaced inputs. Fill in the difference rows,
        then name the function type.

        \vspace{3pt}
        \begin{center}
        \renewcommand{\arraystretch}{1.35}
        \begin{tabular}{|l|c|c|c|c|c|}
        \hline
        $x$    & $0$ & $1$ & $2$ & $3$ & $4$ \\ \hline
        $g(x)$ & $2$ & $5$ & $12$ & $23$ & $38$ \\ \hline
        1st differences & --- & \blank{1cm} & \blank{1cm} & \blank{1cm} & \blank{1cm} \\ \hline
        2nd differences & --- & --- & \blank{1cm} & \blank{1cm} & \blank{1cm} \\ \hline
        \end{tabular}
        \end{center}

        \vspace{2pt}
        Function type: \blank{3cm} \qquad Because \blank{6.4cm}

  \item The parent function is $f(x) = \sqrt{x}$, with domain $[0,\infty)$ and
        range $[0,\infty)$. Let $g(x) = -2f(x - 1) + 4$.

    \begin{enumerate}[label=(\alph*), itemsep=3pt]
      \item Write $g$ in terms of $x$: \quad $g(x) = $ \blank{4.4cm}
      \item The change \emph{inside} the parentheses moves the graph
            \blank{3.4cm}
      \item The changes \emph{outside}, in order: \blank{7cm}
      \item Domain of $g$: \blank{2.6cm} \qquad Range of $g$: \blank{2.6cm}
      \item The parent's landmark point $(0,0)$ lands at \blank{2.4cm} on $g$.
    \end{enumerate}

  \item Two one-to-one functions are given by the table below.

        \vspace{2pt}
        \begin{center}
        \renewcommand{\arraystretch}{1.3}
        \begin{tabular}{|c|c|c|c|c|c|}
        \hline
        $x$    & $0$ & $1$ & $2$ & $3$ & $4$ \\ \hline
        $f(x)$ & $5$ & $2$ & $7$ & $1$ & $3$ \\ \hline
        $g(x)$ & $3$ & $0$ & $4$ & $2$ & $1$ \\ \hline
        \end{tabular}
        \end{center}

    \begin{enumerate}[label=(\alph*), itemsep=3pt]
      \item The handoff in $f(g(1))$ is \blank{1.4cm}, so $f(g(1)) = $
            \blank{1.4cm}
      \item $g(f(3)) = $ \blank{1.4cm}
      \item $f^{-1}(7) = $ \blank{1.4cm} \quad because $f($ \blank{1.1cm} $) = 7$.
      \item How can you tell from the table that $f$ is one-to-one?

            \writeline
      \item Why does $g(f(0))$ have no value from this table?

            \writeline
    \end{enumerate}

  \item Let $f(x) = 3x + 1$ and $g(x) = x^2 - 2$.

\begin{work}
            g(2) &= 4 - 2 \\
                 &= 2 \\
         f(g(2)) &= 3(2) + 1 \\
                 &= 7 \\
            f(2) &= 7 \\
         g(f(2)) &= 7^2 - 2 \\
                 &= 47 \\
  (f \circ g)(x) &= 3(x^2 - 2) + 1 \\
                 &= 3x^2 - 5
\end{work}

    \begin{enumerate}[label=(\alph*), itemsep=3pt]
      \item $f(g(2)) = $ \blank{1.8cm} \qquad \textbf{(b)} $g(f(2)) = $
            \blank{1.8cm}
      \setcounter{enumii}{2}
      \item $(f \circ g)(x) = $ \blank{4.4cm}
      \item What do (a) and (b) show about the order of a composition?

            \writeline
    \end{enumerate}

  \item Let $f(x) = \dfrac{x - 5}{2}$.

\begin{work}
                        y &= \frac{x - 5}{2} \\
                        x &= \frac{y - 5}{2} \\
                       2x &= y - 5 \\
                        y &= 2x + 5 \\
  f\!\left(f^{-1}(x)\right) &= \frac{(2x + 5) - 5}{2} \\
                          &= x
\end{work}

    \begin{enumerate}[label=(\alph*), itemsep=3pt]
      \item Find $f^{-1}$ by swapping and solving: \quad $f^{-1}(x) = $
            \blank{3.4cm}
      \item Verify: \quad $f\!\left(f^{-1}(x)\right) = $ \blank{3.4cm}
      \item $f(11) = 3$, so $f^{-1}(3) = $ \blank{1.6cm}
    \end{enumerate}

\end{enumerate}

% ======================================================================
\boxguard[6]
\parthead{Part D --- Extended Response (18 pts, 9 each)}
Explain your reasoning in full sentences.

\begin{enumerate}[leftmargin=1.9em, itemsep=12pt]

  \item A rectangular dog run is built against a barn wall, so only
        \textbf{three} sides need fencing. The builder has $60$ feet of fence.
        Let $x$ be the length of each of the two sides perpendicular to the barn.

\begin{work}
                          A(x) &= x(60 - 2x) \\
                               &= 60x - 2x^2 \\
                          A(5) &= 5(50) \\
                               &= 250 \\
                         A(10) &= 10(40) \\
                               &= 400 \\
  \frac{A(10) - A(5)}{10 - 5} &= \frac{400 - 250}{5} \\
                               &= 30
\end{work}

    \begin{enumerate}[label=(\alph*), itemsep=3pt]
      \item Which quantity is the \emph{question} the model answers --- its
            input? \blank{4.4cm}
      \item The side parallel to the barn has length \blank{2.4cm}, so
            $A(x) = $ \blank{4cm}
      \item What type of function is $A$, and how do you know?

            \writeline
      \item Domain the context forces: \blank{1.4cm} $< x <$ \blank{1.4cm}
            \quad Why?

            \writeline
      \item Average rate of change of $A$ from $x = 5$ to $x = 10$:
            \blank{1.8cm}, with units \blank{4.4cm}.
      \item State one assumption this model rests on.

            \writeline
    \end{enumerate}

  \item Let $f(x) = x^2$ with domain all real numbers.

    \begin{enumerate}[label=(\alph*), itemsep=3pt]
      \item Explain why $f$ has no inverse function on that domain. Give a
            counterexample pair of inputs.

            \writeline
      \item Give a domain restriction that fixes the problem: \blank{3cm}
            \quad Then $f^{-1}(x) = $ \blank{3cm}
      \item Domain of that inverse: \blank{2.6cm} \qquad Range:
            \blank{2.6cm}
      \item Explain how the graph of a function and the graph of its inverse are
            related, and why.

            \writeline

            \writeline
    \end{enumerate}

\end{enumerate}

\end{document}
